\begin{frame}{Limites}
    \metroset{block=fill}

    \begin{block}{Proposição}
        Seja
        \[ \lim_{n \to \infty} \frac{f(n)}{g(n)} = c \ . \]
        Então \pause
            \begin{itemize}
                \item[(i)] se $c > 0$, então $f(n) = \Theta(g(n))$; \pause
                \item[(ii)] se $c = 0$, então $f(n) = O(g(n))$ e $g(n) \not= O(f(n))$; \pause
                \item[(iii)] se $c = \infty$, então $g(n) = O(f(n))$ e $f(n) \not= O(g(n))$; \pause
            \end{itemize}
    \end{block}


    \begin{block}{Exemplo}
            Temos
                \[\lim_{n \to \infty} \frac{4^n}{2^n} = \lim_{n \to \infty} 2^n = \infty \ ,\]
            portanto $4^n \not = O(2^n)$. De forma geral, se $a > b$ são números reais positivos, então $a^n \not=O(b^n)$
    \end{block}

\end{frame}

